\documentclass[10pt]{article}
\usepackage[T1]{fontenc}
\usepackage[a4paper,margin=19mm]{geometry}
\usepackage{amsmath,amssymb,booktabs,array}
\usepackage{enumitem}
\usepackage{xspace}
\usepackage{xcolor}
\usepackage[colorlinks=true,linkcolor=blue!45!black,urlcolor=blue!45!black]{hyperref}
\definecolor{accent}{HTML}{315B7D}
\setlength{\parindent}{0pt}
\setlength{\parskip}{4pt}
\setlist{nosep,leftmargin=*}
\newcommand{\goal}[1]{\medskip\noindent\textcolor{accent}{\textbf{Experimental goal.}} #1\medskip}
\newcommand{\code}[1]{\texttt{#1}}
\newcommand{\Wten}{\ifmmode\mathrm{W10}\else W10\fi}
\pagestyle{plain}

\title{TimeTensors Experiment 2: Sampling and Batch Size}
\author{}
\date{}
\begin{document}\maketitle
\section{Theoretical overview and goal}
Time-series windows are correlated along both user and date axes.  A mini-batch
therefore defines more than a computational unit: it determines which sources
of variation contribute to a gradient.  Random-window sampling approximates
the global empirical risk, date sampling emphasizes cross-sectional diversity,
and individual sampling emphasizes temporal diversity within selected users.

\goal{Measure how user/date diversity and gradient batch size affect convergence
and final generalization under a fixed optimization budget.}

\section{Implementation details}
The modes are \code{random}, \code{dates}, \code{individuals}, and
\code{all}. The exhaustive profile crosses them with batch sizes 64, 256, and 1024. The
exhaustive mode deterministically exposes every accessible window; other modes
draw according to their named axis.  Evaluation sampling remains fixed so only
training exposure changes.

The default study holds six datasets, three $(L,H)$ settings, DLinear, seeds
1--3, learning rate $10^{-5}$, 10,000 optimizer steps, and 1,000-step
validation constant. It compares random sampling at batch sizes 64, 256, and
1024 with individual sampling at batch size 256. The full profile crosses all
four modes with all three batch sizes. The launcher is
\code{02\_sampling.slurm}.

Report final MSE with seed standard deviation, but also retain interval-average
train loss and validation curves.  Because different sampling modes expose
different numbers of unique windows per optimizer step, include the effective
dataset size and number of optimizer updates when interpreting equal-epoch
runs.

\section{Interpretation}
A batch-size effect can be confounded with diversity: a large batch from one
user is not equivalent to a smaller cross-user batch.  Conclusions should
therefore compare batch size within each mode first, then compare the best
stable configuration across modes.
\end{document}
